% Derivation excerpt, retained from the rank-five census project.
% The current modular speedup is justified in completeness.md.
\section{An independent census}\label{sec:data}
\subsection{Commutativity and coordinates}
The following elementary argument also supplies the commutativity needed
for every signed sphericalization search. The unsigned assertion also follows from the proof below.

\begin{lemma}\label{lem:commutativity}
Every rank-five fusion ring is commutative. A rank-five signed Frobenius
based ring admitting a nonzero, dual-equal dimension character is also
commutative.
\end{lemma}
\begin{proof}
The complex algebra of a (possibly signed) Frobenius based ring is
semisimple: the coefficient-of-unit functional gives the positive
inner product making the basis orthonormal and multiplication a
$*$-representation. A noncommutative algebra of dimension five with a
one-dimensional representation must be $\CC\oplus M_2(\CC)$.
Its unique one-dimensional character is Galois invariant, hence rational,
and its values are algebraic integers, hence integers.
For a fusion ring this character is FP; in the signed assertion it is
the given dimension character. In either case its formal codegree is an
integer $D\geq5$, the sum of five nonzero integer squares.
If $f$ denotes the formal codegree of the two-dimensional representation,
the standard identity $\sum_E\dim(E)/f_E=1$ gives
\[
 f=\frac{2D}{D-1}=2+\frac2{D-1}.
\]
Formal codegrees are algebraic integers. The displayed rational number
is not an integer for $D\geq5$, a contradiction.
\end{proof}

Fix the unit $0$ and one of the three dualities
\[
(0,1,2,3,4),\quad(0,1,2,4,3),\quad(0,2,1,4,3).
\]
Write $C_{ijk}=N_{ij}^{k^*}$. Commutativity and reciprocity say that
$C$ is a symmetric tensor, invariant under simultaneous duality.
The nonunit triples therefore give respectively $20,13,10$ variables,
ordered lexicographically by their smallest orbit representative.
All variables range independently in $[0,M]\cap\ZZ$ before
associativity is imposed. The unit and duality determine the other
entries. Commutativity of the five multiplication matrices is equivalent
to associativity: the common unit is cyclic for the span of the matrices,
and commuting matrices acting identically on the unit agree everywhere.
Direct expansion gives $21,12,13$ distinct quadratic equations,
respectively. The supplement regenerates these polynomials from the
axioms; no source census enters their construction.

\subsection{Self-dual rings: reconstruction from a Krylov matrix}
In the self-dual case use the ordered triples
\[
\begin{gathered}111,112,113,114,122,123,124,133,134,144,\\
222,223,224,233,234,244,333,334,344,444\end{gathered}
\]
as coordinates $v_0,\ldots,v_{19}$. By relabelling, arrange that
$v_1$ is the minimum of the twelve mixed repeated-index coefficients,
and $v_1\leq v_2\leq v_3$. The first nontrivial fusion matrix is
\[
N_1=\begin{pmatrix}0&1&0\\1&a&u^{\mathsf T}\\0&u&Q\end{pmatrix},\qquad
u=(v_1,v_2,v_3)^{\mathsf T},\quad
Q=\begin{pmatrix}v_4&v_5&v_6\\v_5&v_7&v_8\\v_6&v_8&v_9\end{pmatrix},
\quad a=v_0.
\]
For each of the nine bounded coordinates of $(u,Q)$ set
\[
K=[u,Qu,Q^2u],\quad s=u^{\mathsf T}u,\quad t=u^{\mathsf T}Qu,
\quad H=Q^2-aQ+uu^{\mathsf T}-I.
\]
For $i=1,2,3$, let $T_i$ be the bottom-right $3\times3$ block
of $N_{i+1}$; thus $(T_i)_{jk}=C_{i+1,j+1,k+1}$.
The vectors $e_i$ below are the three standard coordinate vectors.
Commuting with $N_1$ gives
\begin{align}
T_i u&=He_i,\label{eq:krylov}\\
T_i Qu&=(QH-sQ+u(Qu)^{\mathsf T})e_i,\nonumber\\
T_i Q^2u&=(Q^2H-sQ^2-tQ+(Qu)(Qu)^{\mathsf T}
                       +u(Q^2u)^{\mathsf T})e_i.\nonumber
\end{align}
These follow successively from the lower blocks of the commutator.
If $\det K\ne0$, multiplication by $\operatorname{adj}(K)$ reconstructs
all ten remaining coefficients, affinely in $a$. One integrality
condition is a linear congruence $A+aB\equiv0\pmod{\det K}$.
A gcd calculation gives every bounded solution for $a$, including the
case $B=0$. For each such solution all remaining divisibilities,
coefficient bounds and all $21$ associativity equations are checked.
Thus there is no search over the last ten coefficients in this stratum.

\subsection{Every singular stratum}
The singular cases are essential to completeness. Multiplication by a
ring element is zero if it annihilates the unit. Hence any polynomial
annihilating the unit under $N_1$ annihilates $N_1$ itself. Apply this to
the cyclic block generated by the unit.

If $\operatorname{rank}K=2$, write $Q^2u=\alpha u+\beta Qu$ and
$\lambda=\operatorname{tr}Q-\beta$. The invariant rational subspaces of
$Q$ show that $\alpha,\beta,\lambda$ are integers. The remaining
one-dimensional eigenspace of $Q$ has eigenvalue $\lambda$. The cyclic
block condition is
\begin{equation}\label{eq:ranktwo}
(\lambda^2-a\lambda-1)(\lambda^2-\beta\lambda-\alpha)
-\lambda(s\lambda+t-\beta s)=0.
\end{equation}
This is linear in $a$. All zero-coefficient cases are retained.

If $\operatorname{rank}K=1$, write $Qu=\lambda u$ and let
$x^2-Tx+U$ be the characteristic polynomial on $u^\perp$.
The cyclic block polynomial is
\[
p(x)=x^3-(a+\lambda)x^2+(a\lambda-1-s)x+\lambda.
\]
For two distinct complementary eigenvalues require $x^2-Tx+U\mid p$;
for a repeated eigenvalue $\eta=T/2$ require $p(\eta)=0$.
These are again linear conditions on $a$, with no division by a
potentially zero coefficient.

For both positive-rank singular cases, the $15$ commutator equations
linear in the ten remaining coordinates are solved by fraction-free
elimination. Enumerate every bounded integral value of the free
coordinates, back-substitute, and test all $21$ quadratic equations.

If $u=0$, the unit is annihilated by $N_1^2-aN_1-I$.
For $a>0$, $x^2-ax-1$ is irreducible over $\Q$, and a five-dimensional
rational vector space cannot be a module over the resulting quadratic
field. Thus $a=0$ and $Q^2=I$. A nonnegative integral symmetric
orthogonal matrix is a permutation matrix. The transpositions are
handled by the same linear completion. For $Q=I$, the remaining ten
variables form the weighted rank-four associativity system: the constant
$1$ in each diagonal rank-four equation is replaced by $2$, because
both $1$ and the invertible element contribute. The Diophantine
elimination of the accompanying rank-four derivation applies with this substitution.
All zero branches are implemented explicitly in \texttt{weighted.hpp}.
For example, its terminal branch $b=c=e=0$ gives $a=1$, $d=f=2$ and
$gi+hj=h^2+i^2-6$; it must not be omitted.

\subsection{Nonself-dual rings and isomorphisms}
For one dual pair, use the $13$ orbit coordinates $v_0,\ldots,v_{12}$.
Put $A=v_6-v_5$, $B=v_{10}-v_9$, $C=v_{11}-v_{12}$.
Three associativity consequences are
\begin{align*}
Cv_5&=Av_2+Bv_4,& Cv_9&=Av_4+Bv_8,\\
v_{11}^2&=1+v_{12}^2-v_5^2+v_6^2-v_9^2+v_{10}^2.
\end{align*}
Enumerate $v_5,v_6,v_9,v_{10},v_{12}$ and recover $v_{11}$ by an exact
square test. Enumerate $v_4$ and solve for $v_2,v_8$, retaining all free
choices if $A$ or $B$ vanishes. Linear commutators determine or bound
the four remaining variables $v_0,v_1,v_3,v_7$; check all equations.
For two dual pairs the corresponding initial reductions are
\[
v_0^2=1+v_1^2-v_2^2-v_3^2+2v_4^2,\qquad
v_7^2=v_2^2-v_3^2+v_5^2.
\]
Two linear equations determine $v_6$ unless their coefficients vanish;
that branch retains every bounded value. Solve for $v_8,v_9$ by linear
completion and check all $13$ equations. The orbit definitions in
\texttt{census\_system2.json} and \texttt{census\_system4.json} fix these
coordinate conventions exactly.

Finally minimize over the permutations preserving the fixed duality:
$24,4,8$ permutations in the three cases. Distinct duality types cannot
be isomorphic. The independent verifier additionally minimizes every
output tensor over all $24$ unit-fixing permutations.

